\section{Executable Specification and Method}
\label{sec:method}

The artifact is a single .NET solution written in F\#.  Purity is enforced at
the transition boundary, not imposed on the surrounding program.  The source
contains ordinary file I/O, JSON parsing, hashing, mutable library internals,
and command-line orchestration outside the reducer.  There is no database,
dependency-injection container, actor framework, plugin system, network client,
or asynchronous interpreter inside the kernel.

\subsection{Artifact structure}

The artifact includes:

\begin{itemize}[leftmargin=*]
  \item a closed event and effect model;
  \item a pure projector and explicit settlement scheduler;
  \item branch construction with parent/child identity, cut, shared prefix,
  continuation, inherited effects, and prefix hash;
  \item cut-level fork assessments and replay-grade computation;
  \item an ActiveGraph/bridge JSONL normalizer and batch validator;
  \item 50 passing tests, including FsCheck properties and named fixtures;
  \item seven language-neutral JSON conformance vectors and a JSON Schema;
  \item read-only SQLite export and actual-fork audit scripts;
  \item a findings log that records failed checked properties as they are
  discovered.
\end{itemize}

All F\# projects treat warnings as errors.  Pattern matches over closed domain
unions are exhaustive.  External effect execution is not linked into the
kernel, so replay cannot accidentally call a model or tool through an injected
client.

\subsection{Property-based checks}

FsCheck generates bounded event values, writer families, effect dimensions,
and step bounds.  Scheduler policies are enumerated rather than sampled because
the executable model exposes four intended extremes.  Properties include:

\begin{itemize}[leftmargin=*]
  \item deterministic projection and event-identity idempotence;
  \item exact-bound quiescence and bounded self-cycle divergence;
  \item convergence of isolated disjoint writers under all schedulers;
  \item schedule dependence for generated conflicting writers;
  \item identity, nested-collapse, projection, and suffix identities for forks;
  \item verdict matrices across generated footprint and replay-source values;
  \item grade-chain mechanics and prerequisite enforcement.
\end{itemize}

The 50-test suite comprises 15 live FsCheck properties and 35 xUnit facts.
Per-property \code{MaxTest} values range from 100 to 250, totaling 3,300
top-level generated trials in a default full run; properties that enumerate
schedulers or dimension values perform additional checks inside each trial.
No fixed global replay seed is configured.  Regression reproducibility comes
from the checked-in named fixtures rather than from repeating one pseudorandom
stream.

FsCheck shrinking is useful for scalar values, but the semantically important
counterexamples are named and checked in as JSON-backed fixtures.  The fixtures
include overlapping writers, read/trigger interference despite disjoint writes,
a self-cycle, a discarded OneShot effect, an unresolved request, malformed
effect transitions, and forged look-alike evidence types.

Table~\ref{tab:vector-coverage} states the deliberately bounded role of the
seven language-neutral vectors.  They are portability smoke tests, not an
exhaustive cross-product of property, footprint, replay source, lifecycle, and
obligation.

\begin{table}[ht]
\centering
\scriptsize
\begin{tabularx}{\textwidth}{@{}P{0.38\textwidth}Y@{}}
\toprule
Vector & Boundary exercised \\
\midrule
\path{projection-last-write-wins} & ProjectionReplay \code{Sound}; deterministic fact projection \\
\path{discarded-one-shot-is-not-counterfactual} & CounterfactualWorld \code{Unsound}; discarded committed OneShot \\
\path{recorded-effect-strict-replay} & StrictExecutionReplay \code{Sound}; committed Recorded result \\
\shortstack[l]{\code{recorded-unknown-footprint-}\\\code{conditions-external-continuation}} & ExternalContinuation \code{Conditional}; unknown inherited footprint \\
\path{outcome-without-request-is-malformed} & ProjectionReplay \code{Unsound}; lifecycle-integrity blocker \\
\path{boundary-grade} & Boundary grade prerequisites \\
\path{native-grade} & Native grade and prerequisite chain \\
\bottomrule
\end{tabularx}
\caption{Coverage intent of the language-neutral v1 conformance vectors.}
\label{tab:vector-coverage}
\end{table}

\subsection{Real-log normalization}

Validation accepts ordered JSONL.  For SQLite-backed ActiveGraph stores, an
export script uses \code{events.seq} as the ordering authority and emits one
file per run.  Source domain events that are not part of the closed effect or
evidence vocabulary remain signals and are reported by type.  The adapter does
not infer effects from every \code{*.requested} event because real domains use
names such as \code{round.requested}; only a closed list of known external
request types is classified.

This choice creates an explicit adapter-coverage premise.  Unknown footprint
fails closed once an event is classified as an effect, but an unclassified
source event is not automatically treated as an effect.  A Sound continuation
verdict therefore means ``Sound under the profile declaration that the listed
unclassified types are domain-only.''  A new source type that hides an external
interaction can invalidate that verdict.  The validator reports every such type
so the declaration can be audited rather than silently assumed.

For an external request, footprint is read from explicit side-effect metadata.
A write without an idempotency key is OneShot; an unknown tool footprint remains
Unknown.  Model and embedding requests default conservatively to OneShot.  A
response is Recorded only when the log contains an explicit response hash,
output hash, or raw output from which the validator derives a content SHA-256.
The derived hash certifies captured replay material; it is not a provider
receipt or proof of external commitment.

Every source event is assigned a contiguous normalized sequence.  Bridge
events can produce multiple evidence facts; this is why the bridge normalized
count exceeds its source count.  The source event identity is retained for the
first fact, and deterministic derived identifiers are used for additional
facts.  Duplicate or non-contiguous normalized identity remains a blocker.  The
normalizer also retains the cumulative cut after each atomic source event.
These source-boundary positions are reported separately from diagnostic cuts
between co-derived facts.

\subsection{All-cut enumeration}

For a normalized log of length $n$, the validator evaluates cuts
$0,\ldots,n$.  Prefix states are computed with one scan; environmental findings
for suffixes are accumulated in reverse.  It reports both (i) every normalized
position, useful for checking the executable transformation, and (ii) only
source-event boundaries, which are the physically meaningful fork candidates
in the source trace.  The validator reports Sound, Conditional, and Unsound
totals for ProjectionReplay, ExternalContinuation, and CounterfactualWorld,
plus the number of runs containing an unsound counterfactual cut.

Actual ActiveGraph forks require a second audit because the child trace has
lineage and store-level sequence differences.  The SQL audit resolves the
parent cut event, orders the parent prefix and child trace independently, and
compares corresponding event ID, type, actor, payload, frame, cause, and
timestamp.  In this SQLite schema, \code{events.seq} is the database-wide
autoincrement primary key, not a run-local envelope field.  It is used to order
each trace and excluded from equality because copied child events receive new
store positions.  The normalizer then assigns the contiguous per-run sequence
used by Equation~\ref{eq:event-envelope}; that derived ordinal is not a second physical field in the
fork audit.  Child run ID is also excluded.  The audit separately counts
retained external requests and unresolved boundaries using the same explicit
request/outcome vocabulary as the normalizer; unknown event types remain
unclassified rather than being inferred as effects.

\subsection{Reproducibility discipline}

Evidence entries record repository, immutable revision when available, source
path, artifact hashes, counts, parser profile, command, publication status, and
limitations.  Checked-in fixtures are hash-bound and explicitly marked
synthetic or sanitized.  Private and local runtime payloads are not copied into
the repository merely to make a test convenient.  Missing revision or source
provenance is reported as a limitation rather than replaced with a hash-only
claim of reproducibility.
